\chapter{Kwantummechanisch impulsmoment}\label{ch:b3:quantum-angular-momentum}

Richt een radiotelescoop op een donkere plek aan de hemel en stem hem af
op \qty{115}{GHz}: er verschijnt een heldere, vlijmscherpe lijn ---
moleculen koolstofmonoxide, tien kelvin boven het absolute nulpunt, die
één trede afdalen op een ladder van \emph{rotatie}toestanden. Hun tuimelen
is, zoals alles wat in de kwantummechanica draait, tweevoudig
gekwantiseerd: de grootte van het impulsmoment kan alleen
$\sqrt{l(l+1)}\,\hbar$ zijn, en zijn component langs een gekozen as alleen
$m\hbar$. Dit hoofdstuk leidt die dubbele kwantisering af --- eenmaal
algebraïsch, uit de \omterm{def:b3:quantum-formalism:commutator}{commutatoren} die het van de \omterm{def:b3:hamiltonian-mechanics:poisson}{poissonhaakjes} van
\cref{ch:b3:hamiltonian-mechanics} erfde, en eenmaal concreet, als de
\omterm{prop:b3:quantum-angular-momentum:harmonics}{bolfuncties} die elke atomaire orbitaal vormgeven. De algebra geeft ons
meer dan we vragen: zij laat halftallige waarden toe die geen enkele
baanbeweging kan verwezenlijken, een vacature die de natuur twee
hoofdstukken verderop met de spin vervult. Onderweg ontmoeten we de
rotatieladders van de moleculen --- de millimeterlijnen waarmee astronomen
het koude gas van de sterrenstelsels wegen --- en de magnetische momenten
waarmee het impulsmoment zich voor het eerst gesplitst liet zien.

\section{De algebra van de rotatie}

\begin{definition}[Operatoren van het impulsmoment]\label{def:b3:quantum-angular-momentum:operators}
\index{impulsmoment!kwantummechanisch}
Het baanimpulsmoment is de operator $\hat{\vect L} =
\hat{\vect r}\wedge\hat{\vect p}$, componentsgewijs $\hat L_z = \hat
x\hat p_y - \hat y\hat p_x$ en cyclisch. Uit $[\hat x, \hat p_x] =
\iu\hbar$ volgt
\[
[\hat L_x, \hat L_y] = \iu\hbar\,\hat L_z
\quad\text{(en cyclisch)} , \qquad
[\hat L^2, \hat L_z] = 0 :
\]
de componenten zijn onderling onverenigbaar --- geen toestand heeft er
twee scherp --- maar het totale kwadraat is met elk ervan verenigbaar:
het paar $(\hat L^2, \hat L_z)$ is de standaardkeuze van etiketten. Deze
\omterm{def:b3:quantum-formalism:commutator}{commutatoren} zijn precies $\iu\hbar$ maal de \omterm{def:b3:hamiltonian-mechanics:poisson}{poissonhaakjes} die in
\cref{ex:b3:hamiltonian-mechanics:brackets} zijn berekend: het woordenboek
van Dirac aan het werk.
\end{definition}

\begin{theorem}[Het spectrum, uit de algebra alleen]\label{thm:b3:quantum-angular-momentum:spectrum}
\index{magnetisch kwantumgetal}
Laat $\hat J_x, \hat J_y, \hat J_z$ \emph{willekeurige} drie \omterm{def:b3:quantum-formalism:observable}{hermitische
operatoren} zijn die aan de bovenstaande commutatiebetrekkingen voldoen.
Dan zijn de gezamenlijke \omterm{def:b3:quantum-formalism:observable}{eigenwaarden} van $(\hat J^2, \hat J_z)$
\[
\hat J^2:\ j(j+1)\hbar^2 , \qquad
\hat J_z:\ m\hbar , \quad m = -j, -j+1, \dots, +j ,
\]
waarin $j$ een niet-negatief \emph{geheel of halftallig} getal is: $2j + 1$ waarden van $m$ voor elke $j$. De \omterm{def:b3:harmonic-oscillator:ladder}{ladderoperatoren} $\hat J_\pm = \hat
J_x \pm \iu\hat J_y$ zetten $m$ met $\pm1$ op bij vaste $j$:
\[
\hat J_\pm\ket{j,m} = \hbar\sqrt{j(j+1) - m(m\pm1)}\,\ket{j,m\pm1} .
\]
\end{theorem}

\begin{proof}[Gedeeltelijk bewijs]
$[\hat J_z, \hat J_\pm] = \pm\hbar\hat J_\pm$: werken met $\hat
J_\pm$
verschuift de \omterm{def:b3:quantum-formalism:observable}{eigenwaarde} van $\hat J_z$ met $\pm\hbar$, bij vaste $\hat J^2$ (die met alles commuteert wat uit de $\hat
J_i$ is opgebouwd). De
norm $\|\hat J_\pm\ket{j,m}\|^2 = \hbar^2[j(j+1) -
m(m\pm1)]$ (berekend
uit $\hat J_\mp\hat J_\pm = \hat J^2 - \hat
J_z^2 \mp \hbar\hat J_z$) moet
niet-negatief blijven: de ladder moet dus bovenaan eindigen bij een
$m_{\max}$ met $m_{\max}(m_{\max}+1) =
j(j+1)$, dus $m_{\max} = j$, en
onderaan bij $m_{\min} = -j$. Van $-j$ naar $+j$ klimmen in stappen van
één dwingt $2j$ een niet-negatief geheel getal te zijn. De notatie
$j(j+1)\hbar^2$ voor de \omterm{def:b3:quantum-formalism:observable}{eigenwaarde} is daarmee achteraf gerechtvaardigd.
\end{proof}

\begin{remark}[Een vacature in de catalogus]\label{rem:b3:quantum-angular-momentum:halfinteger}
De algebra staat $j = \tfrac12, \tfrac32, \dots$ toe --- ladders met een
even aantal treden. De baanbeweging, zo tonen we hierna, gebruikt alleen
gehele getallen. De natuur verspilt echter niets: het elektron zelf draagt
$j = \tfrac12$, zonder enige baan erachter
(\cref{ch:b3:spin-two-level}); de hier afgeleide algebra zal ongewijzigd
het atomaire magnetisme, de kernspins en de qubit aandrijven.
\end{remark}

\section{Baanimpulsmoment: de bolfuncties}

\begin{proposition}[Bolfuncties]\label{prop:b3:quantum-angular-momentum:harmonics}
\index{bolfuncties}\index{baankwantumgetal}
In bolcoördinaten is $\hat L_z = -\iu\hbar\,\partial_\varphi$, en de
gezamenlijke eigenfuncties van $(\hat L^2, \hat L_z)$ op de bol zijn de
\emph{\omterm{prop:b3:quantum-angular-momentum:harmonics}{bolfuncties}} $Y_l^m(\theta, \varphi)$:
\[
\hat L^2\,Y_l^m = l(l+1)\hbar^2\,Y_l^m , \qquad
\hat L_z\,Y_l^m = m\hbar\,Y_l^m ,
\]
met $l = 0, 1, 2, \dots$ \emph{geheel} --- de eenwaardigheid van
$\eu^{\iu m\varphi}$ dwingt $m$ geheel te zijn, en dus ook $l$. De eerste
paar, op de normering na: $Y_0^0 = \text{const}$ (de vorm $s$, een bol);
$Y_1^0 \propto \cos\theta$ en $Y_1^{\pm1}
\propto \sin\theta\,\eu^{\pm\iu\varphi}$ (de vormen $p$, twee lobben); $Y_2^0 \propto 3\cos^2\theta - 1$ en haar partners (de familie $d$). Pariteit:
$Y_l^m(-\vect r$ als richting$) = (-1)^l\,Y_l^m$.
\end{proposition}
\admitted

\begin{omfigure}
\begin{tikzpicture}[font=\small, scale=0.94]
  % s orbital
  \begin{scope}
    \draw[->] (0,-1.7) -- (0,1.8);
    \node[anchor=south] at (0,1.8) {$z$};
    \draw[omDef, thick, fill=omDef!18] (0,0) circle (1.1);
    \node at (0,-2.3) {$l = 0$ ($s$): bol};
  \end{scope}
  % p_z
  \begin{scope}[xshift=4.4cm]
    \draw[->] (0,-1.7) -- (0,1.8);
    \node[anchor=south] at (0,1.8) {$z$};
    \draw[omThm, thick, fill=omThm!18, domain=0:360, smooth, samples=100]
      plot ({1.45*abs(cos(\x))*sin(\x)}, {1.45*abs(cos(\x))*cos(\x)});
    \node at (0,-2.3) {$l = 1$, $m = 0$: $|\cos\theta|$};
  \end{scope}
  % p_+-1
  \begin{scope}[xshift=8.8cm]
    \draw[->] (0,-1.7) -- (0,1.8);
    \node[anchor=south] at (0,1.8) {$z$};
    \draw[omExo, thick, fill=omExo!18, domain=0:360, smooth, samples=100]
      plot ({1.45*abs(sin(\x))*sin(\x)}, {1.45*abs(sin(\x))*cos(\x)});
    \node at (0,-2.3) {$l = 1$, $m = \pm1$: $|\sin\theta|$};
  \end{scope}
  % d_z2
  \begin{scope}[xshift=13.2cm]
    \draw[->] (0,-1.7) -- (0,1.8);
    \node[anchor=south] at (0,1.8) {$z$};
    \draw[omProp, thick, fill=omProp!18, domain=0:360, smooth, samples=140]
      plot ({0.62*abs(3*cos(\x)*cos(\x)-1)*sin(\x)}, {0.62*abs(3*cos(\x)*cos(\x)-1)*cos(\x)});
    \node at (0,-2.3) {$l = 2$, $m = 0$};
  \end{scope}
\end{tikzpicture}

{\small Poolgrafieken van $|Y_l^m(\theta)|$ (doorsnede in een vlak door de
$z$-as; de volledige vorm is het omwentelingslichaam). Deze hoekskeletten,
onafhankelijk van welke potentiaal ook, zullen in
\cref{ch:b3:hydrogen-atom} elk atoom aankleden.}
\end{omfigure}

\begin{example}[De starre rotor en zijn ladder]\label{ex:b3:quantum-angular-momentum:rotor}
Een tweeatomig molecuul dat tuimelt met traagheidsmoment $I$ heeft $\hat H
= \hat L^2/2I$: energieën
\[
E_J = \frac{\hbar^2}{2I}\,J(J+1) = B\,J(J+1) ,
\qquad J = 0, 1, 2, \dots
\]
elk $(2J+1)$-voudig ontaard. Fotonabsorptie voldoet aan $\Delta J =
\pm1$,
zodat de absorptiefrequenties $\nu_{J+1\leftarrow J} =
2B(J+1)/h$ zijn:
een kam van \emph{gelijk verdeelde} lijnen --- het kenmerk waaraan men een
rotatiespectrum in één oogopslag herkent. Voor koolstofmonoxide is $B/h = \qty{57.6}{GHz}$: de grondlijn valt bij \qty{115}{GHz} ($\lambda = \qty{2.6}{mm}$), het werkpaard van de millimeterradioastronomie
(\cref{pb:b3:quantum-angular-momentum:1}).
\end{example}

\section{Centrale potentialen, voltooid}

\begin{theorem}[Scheiding in elke centrale potentiaal]\label{thm:b3:quantum-angular-momentum:radial}
\index{radiale vergelijking}\index{centrifugale barrière}
Voor $V(r)$ mogen de stationaire toestanden worden genomen als
\[
\psi(\vect r) = \frac{u(r)}{r}\,Y_l^m(\theta, \varphi) ,
\]
waarin $u$ de eendimensionale \omterm{prop:b3:schrodinger-three-dimensions:swave}{radiale vergelijking} oplost
\[
-\frac{\hbar^2}{2m}\,u'' + \Big[V(r) +
\frac{\hbar^2\,l(l+1)}{2mr^2}\Big]u = E\,u , \qquad u(0) = 0 .
\]
Het hoekvraagstuk is eens en voor altijd door de $Y_l^m$ opgelost; elke
$l$ voegt de afstotende \emph{\omterm{thm:b3:quantum-angular-momentum:radial}{centrifugale barrière}}
$\hbar^2l(l+1)/2mr^2$ toe --- de kwantumversie van de effectieve
potentiaal van \cref{ex:b3:lagrangian-mechanics:central} --- en elk niveau
bij gegeven $l$ is $(2l+1)$-voudig ontaard, want geen centrale kracht kan
zich om de oriëntatie van de $z$-as bekommeren. De $s$-golftruc van
\cref{prop:b3:schrodinger-three-dimensions:swave} was de rij $l = 0$ van
deze stelling.
\end{theorem}

\begin{proof}[Gedeeltelijk bewijs]
De laplaciaan splitst als $\Delta = \tfrac1r\partial_r^2(r\,\cdot) -
\hat L^2/\hbar^2r^2$ (aangenomen; het is de uitspraak dat $\hat L^2$ het
hoekdeel van $-\hbar^2\Delta$ is). Vul $\psi = (u/r)Y_l^m$ in en gebruik
de \omterm{def:b3:quantum-formalism:observable}{eigenwaarde} van $\hat L^2$: de aangegeven vergelijking. De ontaarding
volgt omdat $\hat H$ met alle drie de $\hat
L_i$ commuteert, waarvan de
\omterm{def:b3:harmonic-oscillator:ladder}{ladderoperatoren} $m$ verschuiven zonder de energie te veranderen.
\end{proof}

\begin{omfigure}
\begin{tikzpicture}
  \begin{axis}[omaxis, axis x line=bottom, axis y line=left,
    width=10.6cm, height=6.0cm, xmin=0, xmax=8, ymin=-1.35, ymax=1.3,
    xlabel={$r$ (in eenheden $a_0$)}, ylabel={$U_{\text{eff}}$ (in eenheden $E_{\text{I}}$)},
    xlabel style={at={(axis description cs:0.5,-0.1)}, anchor=north},
    xtick={2,4,6}, ytick=\empty,
    legend style={font=\scriptsize, at={(0.97,0.95)}, anchor=north east,
      draw=none, fill=none}]
    \addplot[omDef, thick, domain=1.35:8, samples=120] {-2/x};
    \addlegendentry{$l = 0$: zuiver coulombs}
    \addplot[omThm, thick, domain=0.42:8, samples=160] {-2/x + 1/(x*x)};
    \addlegendentry{$l = 1$}
    \addplot[omExo, thick, domain=0.85:8, samples=160] {-2/x + 3/(x*x)};
    \addlegendentry{$l = 2$}
    \draw[gray, dashed] (axis cs:0,0) -- (axis cs:8,0);
  \end{axis}
\end{tikzpicture}

{\small Effectieve radiale potentialen van het coulombvraagstuk: de
\omterm{thm:b3:quantum-angular-momentum:radial}{centrifugale barrière} $\hbar^2l(l+1)/2mr^2$ muurt de oorsprong af voor $l \ge 1$. Alleen $s$-toestanden raken de kern --- met gevolgen van de
atoomspectra tot de radioactieve elektronvangst.}
\end{omfigure}

\section{Magnetische momenten: impulsmoment zichtbaar gemaakt}

\begin{proposition}[Magnetisch baanmoment en het zeemaneffect]\label{prop:b3:quantum-angular-momentum:magnetic}
\index{bohrmagneton}\index{zeemaneffect}
Een deeltje met lading $q$ en massa $m$ met baanimpulsmoment
$\hat{\vect L}$ draagt het magnetische moment
\[
\hat{\vect\mu} = \frac{q}{2m}\,\hat{\vect L} ;
\]
voor het elektron is de natuurlijke eenheid het \emph{\omterm{prop:b3:quantum-angular-momentum:magnetic}{bohrmagneton}}
$\mu_{\text{B}} = e\hbar/2m_{\text{e}} = \qty{5.79e-5}{eV/T}$
(vergelijk \cref{exo:b3:schrodinger-three-dimensions:5}). In een veld
$B\vect e_z$ verschuift de energie $-\hat{\vect\mu}\cdot\vect B$ elk
niveau met $+m\,\mu_{\text{B}}B$ (de lading van het elektron is
negatief): een niveau met gegeven $l$ splitst in zijn $2l + 1$
componenten --- het \emph{\omterm{prop:b3:quantum-angular-momentum:magnetic}{zeemaneffect}}, de eerste rechtstreekse
vertoning van de kwantisering van $m$, en de reden dat $m$ het
magnetische kwantumgetal heet. Bij $B = \qty{1}{T}$ bedraagt de
splitsing $\qty{58}{\micro eV}$: klein naast optische energieën, moeiteloos
op te lossen als een verschuiving van spectraallijnen --- en, omgekeerd
gelezen, een magnetometer: de zeemansplitsing van de lijnen van het
zonlicht brengt de magneetvelden van zonnevlekken in kaart.
\end{proposition}

\begin{proof}[Gedeeltelijk bewijs]
Klassiek is een lading op een baan een stroomlus: $\mu = IA =
(qv/2\pi r)(\pi r^2) = qL/2m$; de uitspraak voor operatoren erft dat (en
volgt uit de koppeling aan $\hat{\vect A}$ van
\cref{prop:b3:lagrangian-mechanics:charge}). De energieverschuiving is
storingsrekening in eerste orde, hier vooruitgenomen en in
\cref{ch:b3:perturbation-theory} gerechtvaardigd.
\end{proof}

\begin{omfigure}
\begin{tikzpicture}[font=\small, >=Stealth]
  \draw[->] (0,-3.0) -- (0,3.15);
  \node[anchor=south] at (0,3.15) {$z$};
  % sphere radius sqrt6
  \pgfmathsetmacro\R{2.449}
  \draw[gray, dashed] (0,0) circle (\R);
  \foreach \m/\c in {2/omDef, 1/omThm, 0/omExo, -1/omThm, -2/omDef} {
    \pgfmathsetmacro\zz{\m}
    \pgfmathsetmacro\rr{sqrt(\R*\R-\m*\m)}
    \draw[\c, thick, ->] (0,0) -- (\rr,\zz);
    \draw[gray!60, thin, dashed] (-0.05,\zz) -- (\rr,\zz);
    \node[anchor=west, \c] at ({\rr+0.1},\zz) {$m = \m$};
  }
  \node[anchor=east, font=\footnotesize] at (-0.15,2.0) {$L_z = m\hbar$};
  \node[anchor=north, font=\footnotesize] at (1.3,-2.75) {$\|\vect L\| = \sqrt{6}\,\hbar$};
\end{tikzpicture}

{\small ``Ruimtekwantisering'' voor $l = 2$: de vector van het
impulsmoment heeft lengte $\sqrt{l(l+1)}\,\hbar = \sqrt6\,\hbar$ maar kan
de $z$-as slechts de vijf projecties $m\hbar$ aanbieden --- nooit haar
volle lengte: omdat de componenten onverenigbaar zijn, kan $\vect L$ nooit
precies langs een as liggen.}
\end{omfigure}

\begin{method}[Impulsmoment in de praktijk]\label{met:b3:quantum-angular-momentum:method}
(1) Etiketteer toestanden met $(l, m)$ --- of met $(j, m)$ voor de
abstracte algebra --- en vraag nooit naar twee componenten tegelijk. (2)
Matrixelementen: gebruik de ladderformules; $\hat L_x = (\hat L_+ + \hat
L_-)/2$. (3) Centrale potentiaal: haal de $Y_l^m$ aan en los alleen de
\omterm{prop:b3:schrodinger-three-dimensions:swave}{radiale vergelijking} met de \omterm{thm:b3:quantum-angular-momentum:radial}{centrifugale barrière} op. (4)
Rotatie-energieën: $BJ(J+1)$, lijnen bij $2B(J+1)$ --- haal $B$ en dus
bindingslengten uit de gelijke afstanden. (5) Magnetische vragen: zet
impulsmomenten om in momenten met $\mu_{\text{B}}$ per $\hbar$, en
energieën met $\mu_{\text{B}}B$ per eenheid $m$.
\end{method}

\section{Opgaven}

\begin{exercise}[$\star$]\label{exo:b3:quantum-angular-momentum:1}
(a) Leid $[\hat L_x, \hat L_y] = \iu\hbar\hat L_z$ af uit de canonieke
\omterm{def:b3:quantum-formalism:commutator}{commutatoren}. (b) Toon aan dat $[\hat L^2, \hat L_z] = 0$. (c) Toon aan
dat $[\hat L_z, \hat x] = \iu\hbar\hat y$: wat brengt $\hat L_z$ voort?
(d) Waarom laten de drie componenten geen gemeenschappelijke eigenbasis
toe --- op één bijzondere toestand na (welke)?
\end{exercise}

\begin{exercise}[$\star$]\label{exo:b3:quantum-angular-momentum:2}
Voor $l = 1$, in de basis $\{\ket{1,1}, \ket{1,0}, \ket{1,-1}\}$:
(a) schrijf de matrix van $\hat L_z$ op; (b) schrijf met de ladderformule
$\hat L_+$ en $\hat L_-$ op; (c) stel $\hat L_x$ samen en ga na dat haar
\omterm{def:b3:quantum-formalism:observable}{eigenwaarden} $\hbar, 0, -\hbar$ zijn; (d) een toestand met $L_x =
+\hbar$
wordt langs $z$ gemeten: geef de kansen op de drie uitkomsten.
\end{exercise}

\begin{exercise}[$\star$]\label{exo:b3:quantum-angular-momentum:3}
Koolstofmonoxide: bindingslengte \qty{0.113}{nm}, gereduceerde massa
$\qty{1.14e-26}{kg}$. (a) Bereken $I$ en $B = \hbar^2/2I$ in meV.
(b) De frequentie en de golflengte van de lijn $J = 1 \leftarrow 0$.
(c) De volgende twee lijnen. (d) Een astronome meet de afstand in de kam
op vijf cijfers: welke moleculaire grootheid krijgt zij daarmee in
handen, en met welke precisie?
\end{exercise}

\begin{exercise}[$\star$]\label{exo:b3:quantum-angular-momentum:4}
Een deeltje bevindt zich in een toestand met $l = 2$. (a) Geef de mogelijke
uitkomsten van een meting van $L_z$. (b) De kleinste hoek tussen $\vect L$
en de $z$-as, met $\cos\theta = m/\sqrt{l(l+1)}$: reken haar uit. (c)
Waarom kan de hoek nooit nul zijn? (d) Toon aan dat de minimale hoek naar
nul gaat als $l \to \infty$: de klassieke vectoren teruggevonden.
\end{exercise}

\begin{exercise}[$\star\star$]\label{exo:b3:quantum-angular-momentum:5}
Het hoekdeel van $-\hbar^2\Delta$ is $\hat L^2$, met
\[
\hat L^2 = -\hbar^2\Big[\frac{1}{\sin\theta}\,\partial_\theta
(\sin\theta\,\partial_\theta) +
\frac{1}{\sin^2\theta}\,\partial_\varphi^2\Big] .
\]
(a) Ga na dat $Y_1^0 \propto \cos\theta$ de $\hat L^2$-eigenwaarde
$2\hbar^2$ heeft. (b) Ga hetzelfde na voor $Y_1^{\pm1} \propto \sin\theta\,
\eu^{\pm\iu\varphi}$, met hun \omterm{def:b3:quantum-formalism:observable}{eigenwaarden} van $\hat L_z$. (c)
Controleer de pariteiten. (d) Waarom moet $\braket{Y_1^0}{Y_1^1} = 0$ zijn
zonder ook maar één integraal uit te rekenen?
\end{exercise}

\begin{exercise}[$\star\star$]\label{exo:b3:quantum-angular-momentum:6}
Welke rotatielijn is de helderste? De bezetting van niveau $J$ is
$\propto (2J+1)\,\eu^{-BJ(J+1)/k_{\text{B}}T}$. (a) Verklaar de twee
factoren. (b) Toon aan dat het maximum nabij $J_{\max} \approx
\sqrt{k_{\text{B}}T/2B} - \tfrac12$ ligt. (c) Voor CO bij \qty{10}{K}
(een donkere wolk) en bij \qty{300}{K}: welke lijnen overheersen? (d)
Omgekeerd: welke temperatuur verraadt een waargenomen CO-ladder die bij
$J = 7$ piekt?
\end{exercise}

\begin{exercise}[$\star\star$]\label{exo:b3:quantum-angular-momentum:7}
Normaal \omterm{prop:b3:quantum-angular-momentum:magnetic}{zeemaneffect}. Een spectraallijn bij \qty{500}{nm} komt van een
overgang $l = 2 \to l = 1$ in een veld $B = \qty{2}{T}$; verwaarloos de
spin (geldig voor bijzondere ``singlet''toestanden). (a) Schets de
gesplitste niveaus. (b) Toon met de selectieregel $\Delta m = 0, \pm1$ aan
dat er slechts \emph{drie} lijnposities verschijnen, bij $0, \pm\mu_{\text{B}}B/h$. (c) Bereken de splitsing in GHz en in picometers
golflengte. (d) De meeste echte lijnen splitsen in meer dan drie
componenten (het ``anomale'' \omterm{prop:b3:quantum-angular-momentum:magnetic}{zeemaneffect}): van welk ontbrekend
ingrediënt, door het elektron zelf gedragen, was dat historisch het bewijs?
\end{exercise}

\begin{exercise}[$\star\star$]\label{exo:b3:quantum-angular-momentum:8}
De centrifugale wand. Schrijf voor waterstof ($V = -e^2/4\pi\varepsilon_0r$) de $U_{\text{eff}}$ voor $l = 1$ op in eenheden
$E_{\text{I}}$ en $a_0$. (a) Bepaal het minimum en zijn waarde. (b) Toon
aan dat $U_{\text{eff}} >
0$ voor $r < \dots$ --- bepaal de straal
waarbinnen de barrière overheerst. (c) Vergelijk $|\psi|^2$ nabij $r = 0$
voor $s$- en $p$-toestanden ($u \sim r^{l+1}$, aangenomen): welke
toestanden ``raken'' de kern? (d) Elektronvangst --- een kern die een van
de elektronen van zijn atoom opslokt --- verloopt overweldigend vanuit
$s$-schillen: leg dat in één zin uit.
\end{exercise}

\begin{exercise}[$\star\star$]\label{exo:b3:quantum-angular-momentum:9}
Op $\ket{l, m}$: (a) toon aan dat $\langle\hat L_x\rangle = \langle\hat
L_y\rangle = 0$; (b) toon aan dat $\langle\hat L_x^2\rangle = \langle\hat
L_y^2\rangle = \tfrac12[l(l+1) - m^2]\hbar^2$; (c) ga de
\omterm{thm:b3:quantum-formalism:uncertainty}{onzekerheidsrelatie} $\Delta L_x\,\Delta L_y \ge \tfrac\hbar2
|\langle\hat L_z\rangle|$ na op de uitgerekte toestand $m = l$; (d) duid
het ``kegel''beeld van de figuur in het licht van (a) en (b).
\end{exercise}

\begin{exercise}[$\star\star\star$]\label{exo:b3:quantum-angular-momentum:10}
Trillings-rotatiebanden. Een infrarode trillingsovergang ($\hbar\omega_0$)
van een tweeatomig molecuul verandert tegelijk $J$ met $\pm1$. (a) Toon
aan dat de absorptielijnen vallen bij $h\nu =
\hbar\omega_0 + 2B(J+1)$ (de
tak $R$, van $J \to J+1$) en $h\nu = \hbar\omega_0 - 2BJ$ (de tak $P$, $J \to J-1$), met $J \ge
1$. (b) Waarom is er een \emph{gat} bij
$\hbar\omega_0$ zelf? (c) Schets de band: twee kammen die een ontbrekende
middellijn flankeren, waarbij de sterkte van elke lijn
\cref{exo:b3:quantum-angular-momentum:6} volgt. (d) De band bij
\qty{15}{\micro m} van \cref{pb:b3:harmonic-oscillator:1} heeft precies
deze anatomie: wat bepaalt haar totale breedte, en dus de ``flanken'' die
extra koolstofdioxide werkzaam maken?
\end{exercise}

\begin{exercise}[$\star\star\star$]\label{exo:b3:quantum-angular-momentum:11}
Maak de algebra af. (a) Bereken uit $\hat J_\mp\hat J_\pm = \hat J^2 -
\hat J_z^2 \mp \hbar\hat J_z$ de waarde $\|\hat J_\pm\ket{j,m}\|^2$ en de
laddercoëfficiënten. (b) Toon aan dat, als de ladder niet zou eindigen,
een norm negatief zou worden: wijs de schuldige toestand aan. (c)
Concludeer dat $m_{\max} - m_{\min} = 2j$ een niet-negatief geheel getal
moet zijn en som de toegestane $j$ op. (d) Waar precies \emph{faalt} het
argument om halftallige waarden uit te sluiten --- en welke bijkomende eis
(eenwaardigheid op de bol) sluit ze alleen voor de baanbeweging uit?
\end{exercise}

\begin{exercise}[$\star\star\star$]\label{exo:b3:quantum-angular-momentum:12}
Einstein--De Haas: impulsmoment is mechanisch. Een ijzeren cilinder
(straal $a = \qty{1.0}{cm}$, dichtheid $\rho = \qty{7.9e3}{kg/m^3}$, $n = \qty{8.5e28}{atoms/m^3}$) hangt aan een torsiedraad; het omkeren van zijn
magnetisatie keert ongeveer één $\hbar$ aan impulsmoment per atoom om. (a)
Volgens het behoud moet de staaf terugstoten: bereken het overgedragen
impulsmoment per eenheid van volume. (b) Toon met het traagheidsmoment van de staaf, $\tfrac12 M
a^2$, aan dat
de staaf $\omega = 2n\hbar/\rho a^2$ krijgt en reken dit uit. (c) De proef
van 1915 dreef de omkering op de resonantie van de draad aan om de
minuscule schoppen op te tellen: schat de trillingsamplitude na $Q \sim 100$ resonante omkeringen, uitgaande van uw $\omega$ per schop. (d) De
gemeten verhouding van moment tot impulsmoment kwam uit op tweemaal de
baanvoorspelling $e/2m_{\text{e}}$: wat kondigde die factor twee aan?
\end{exercise}

\section{Vraagstuk: Sterrenstelsels wegen op 2,6 millimeter}

\begin{problem}[{Weekendvraagstuk --- koolstofmonoxide, radioastronomie
en het koude heelal}]\label{pb:b3:quantum-angular-momentum:1}
Het meeste stervormende materiaal van een sterrenstelsel is koude
moleculaire waterstof --- en die is onzichtbaar: H$_2$ is symmetrisch,
heeft geen dipoolmoment en geen rotatielijnen bij de temperaturen van
koude wolken. De oplossing van de astronomie is haar trouwe metgezel. Eén
CO-molecuul per $\num{e4}$ H$_2$ zet met zijn ladder bij \qty{115}{GHz}
elke moleculaire wolk aan de hemel aan. Gegevens: voor CO is $B/h = \qty{57.6}{GHz}$ ($B =
\qty{0.238}{meV}$); bindingslengte $r_0 = \qty{0.113}{nm}$; gereduceerde massa $\mu = \qty{1.14e-26}{kg}$;
$k_{\text{B}}T$ bij \qty{10}{K} is \qty{0.862}{meV}.

\medskip\noindent\textbf{Deel I --- De kwantumrotor.}
\begin{enumerate}
  \item Geef uit $\hat H = \hat L^2/2I$ de energieën $E_J$ en hun
        ontaardingen.
  \item Ga $B = \hbar^2/2I$ voor CO na uit $r_0$ en $\mu$.
  \item Leid met de selectieregel $\Delta J = \pm1$ de
        absorptiefrequenties $\nu_{J+1\leftarrow J} = 2B(J+1)/h$ af en
        noem de eerste drie.
  \item Waarom liggen de lijnen \emph{gelijk verdeeld} --- en wat geeft
        een gemeten afstand de astronoom in handen (via $I$)?
  \item Bereken de golflengte van de grondlijn en leg uit waarom het
        waarnemen ervan hoge, droge plekken vergt (wat absorbeert
        millimetergolven in onze eigen atmosfeer --- denk aan de buur van
        het molecuul uit \cref{pb:b3:harmonic-oscillator:1}, het water)?
  \item Het foton van de lijn $1 \to 0$ draagt $\hbar$ aan impulsmoment:
        waarom is $\Delta J = \pm1$ niet slechts een vuistregel maar een
        behoudswet?
\end{enumerate}

\medskip\noindent\textbf{Deel II --- Waarom CO en niet H$_2$.}
\begin{enumerate}[resume]
  \item H$_2$ is homonucleair: zeg waarom het in het geheel geen
        rotatiedipoolstraling uitzendt.
  \item H$_2$ is bovendien \emph{licht}: bereken de rotatieconstante van
        H$_2$ ($\mu = m_{\text{p}}/2$, $r_0 = \qty{0.074}{nm}$) en de
        energie van zijn eerste bereikbare overgang; vergelijk met
        $k_{\text{B}}T$ bij \qty{10}{K}: zou koude H$_2$ zelfs via
        quadrupoolachterdeurtjes kunnen stralen?
  \item De eerste trede van CO kost $2B = \qty{0.48}{meV}$ tegen
        $k_{\text{B}}T = \qty{0.86}{meV}$: leeft de ladder bij
        \qty{10}{K}?
  \item Bereken de bezettingsverhouding $n_1/n_0$ bij \qty{10}{K}
        (ontaarding meegerekend).
  \item Nabij welke $J$ piekt de bezetting bij \qty{10}{K} (gebruik
        $J_{\max} \approx \sqrt{k_{\text{B}}T/2B} - \tfrac12$)?
  \item Vat in één zin het pact van de ruil samen: wat CO levert, en wat
        er over de verhouding CO tot H$_2$ moet worden aangenomen.
\end{enumerate}

\medskip\noindent\textbf{Deel III --- De hemel lezen.}
\begin{enumerate}[resume]
  \item De lijn van een wolk komt aan op \qty{115.156}{GHz} in plaats van
        de laboratoriumwaarde \qty{115.271}{GHz}: bereken de snelheid van
        de wolk langs de gezichtslijn, en haar richting.
  \item De lijn wordt door inwendige bewegingen van $\pm\qty{2}{km/s}$
        doppler\emph{verbreed}: bereken de lijnbreedte in MHz en
        vergelijk met de thermische breedte die bij \qty{10}{K} wordt
        verwacht ($v_{\text{th}} \sim \sqrt{k_{\text{B}}T/m_{\text{CO}}}
        \approx \qty{55}{m/s}$): wat
        overheerst de verbreding?
  \item De dopplerverschuiving van CO over de schijf van een
        spiraalstelsel in kaart brengen levert zijn rotatiekromme
        $v(r)$. Die krommen blijven ver buiten de zichtbare schijf
        \emph{vlak}: haal op wat $v(r)$ rond een centraal
        geconcentreerde massa zou moeten doen (volume van jaar 1,
        gravitatie), en zeg wat de vlakheid impliceert.
  \item De verhouding van de helderheid van de lijn $2\to1$ tot die van
        $1\to0$ meet de bezettingen van de niveaus: welke ene fysische
        grootheid van de wolk levert zij?
  \item ALMA lost CO op in sterrenstelsels waarvan het licht vertrok toen
        het heelal jong was; de waargenomen frequentie van de lijn
        $1 \to 0$ van een stelsel bij ``roodverschuiving $z = 1$'' komt
        aan op de helft van haar rustwaarde --- bij welke frequentie, en
        (denk aan \cref{ex:b3:relativistic-kinematics:redshift}) wat rekte
        haar uit?
  \item Uit een gemeten lijnlichtkracht van CO geven astronomen
        wolkmassa's in zonsmassa's op: som de keten van omrekeningen op
        (lijnfotonen $\to$ aantal CO $\to$ massa H$_2$) en noem de zwakste
        schakel.
\end{enumerate}

\medskip\noindent\textbf{Deel IV --- De koude kraamkamers.}
\begin{enumerate}[resume]
  \item Stervormende kernen liggen op \qty{10}{K}: volgens de wet van
        Wien piekt hun thermische gloed nabij \qty{290}{\micro m} ---
        onzichtbaar voor elke optische telescoop. In één zin: waarom is
        de rotatieladder de \emph{enige} thermometer en maatstaf die zo'n
        kern biedt?
  \item Het niveau $J = 1$ ligt \qty{5.5}{K} (in temperatuureenheden)
        boven de grondtoestand: waarom maakt dat de lijn $1\to0$ tot een
        voortreffelijke thermometer juist in het regime van \qty{10}{K}
        (en niet, zeg, een optische lijn bij \qty{2}{eV} $\sim \qty{23000}{K}$)?
  \item Waarom is een \emph{instortende} kern uiteindelijk niet meer in
        CO ($1\to0$) zichtbaar --- denk aan wat hoge dichtheid en stof met
        de lijn en haar ontsnapping doen.
  \item Moleculaire wolken vertonen ook lijnen van NH$_3$, HCN en
        tientallen andere, elk met haar eigen $B$ en dipoolmoment: wat
        laat de rijkdom van deze ``chemische radioschaal'' astronomen
        ontwarren?
  \item De hele hemel gloeit op \qty{2.7}{K} (de kosmische
        achtergrondstraling): wat gebeurt er met het contrast van de
        CO-lijnen van een wolk naarmate haar eigen temperatuur
        \qty{2.7}{K} nadert, en waarom is dat een harde bodem voor deze
        thermometer?
  \item Een millimeterschotel moet haar vorm tot op ongeveer
        $\lambda/20$ vasthouden: bereken die tolerantie bij
        \qty{115}{GHz} en vergelijk met de tolerantie waaraan een
        optische spiegel (\qty{500}{nm}) moet voldoen --- waarom konden
        de schotels van twaalf meter van ALMA in de open lucht worden
        gebouwd?
  \item Vat het genoemde resultaat samen: een ladder $BJ(J+1)$ met $B/h =
        \qty{57.6}{GHz}$, bezet bij tien kelvin en afgelezen op
        \qty{2.6}{mm}, is hoe de massa, de temperatuur, de snelheid en de
        rotatie van het koude heelal --- en het bewijs voor donkere
        materie rond spiraalstelsels --- worden gemeten.
\end{enumerate}
\end{problem}
